\begin{frame}{Spectral Element Method}
  The defining feature of the spectral element method (SEM) is the selection of internal element nodes corresponding to the Gauss-Lobatto-Legendre (GLL) rule.
  From \cite{devilleHighorderMethodsIncompressible2002}:
  \begin{equation}
    \label{eq:gll}
    \begin{aligned}
       & \int_{-1}^1 u(\xi)\odif{\xi} & \approx & \sum_{k=0}^{N}\rho_k u\left(\xi_k\right)                                  \\
       & \xi_k                        & =       & \begin{cases}
                                                    -1,                        & k=0              \\
                                                    \text{zeros of }L'_N(\xi), & 1\leq k \leq N-1 \\
                                                    1,                         & k=N              \\
                                                  \end{cases}                           \\
       & \rho_k                       & =       & \frac{2}{N\left(N+1\right)}\frac{1}{\left[L_N\left(\xi_k\right)\right]^2}
    \end{aligned}
  \end{equation}
  The above rule is exact for polynomials degree up to $2N-1$.
\end{frame}

\begin{frame}{Spectral Element Method}
  Using the points in \cref{eq:gll}, which we approximate the integrals for the stiffness and mass matrices.
  Using \textit{generic} indices $i,j$, we have:
  \begin{subequations}
    \begin{equation}
      \left(\hat{M}_n\right)_{ij}  \approx \frac{L_n}{2}\sum_{k=0}^{N}\rho_k\ell_i (\mu_k)\ell_j (\mu_k)=\begin{cases}
        \rho_i, & i=j     \\
        0,      & i\neq j
      \end{cases}
    \end{equation}
    \begin{equation}
      \left(\hat{K}_n\right)_{ij}  \approx \frac{2}{L_n}\sum_{k=0}^{N}\rho_k \odv{\ell_i}{\mu}(\mu_k)\odv{\ell_j}{\mu}(\mu_k)
    \end{equation}
  \end{subequations}
  Notice that because the GLL quadrature points are also the nodes for the Lagrange interpolants, the mass matrix is \emph{diagonal}.
  The stiffness matrix, however, is full.
\end{frame}